\documentclass{report}
\usepackage{amsmath,amssymb}
\setlength{\parindent}{0mm}
\setlength{\parskip}{1em}
\begin{document}
\begin{center}
\rule{6in}{1pt} \
{\large Geoffrey Sanders \\
{\bf Eigenvectors of Matrices Associated with Scale-Free Graphs}}

Lawrence Livermore National Laboratory \\ Box 808 \\ L-561 \\ Livermore \\ CA 94551-0808
\\
{\tt sanders29@llnl.gov}\\
Van Emden Henson\\
Panayot Vassilevski\end{center}

A scale-free graph is characterized as having vertices with a degree
distribution that follows a power law. Clustering, ranking, and measuring
similarity of vertices on scale-free graphs are all quantities that
network scientists often calculate. We describe a few network
calculations that are obtained by solving eigensystems or linear systems
involving matrices whose sparsity structure is related to an underlying
scale-free graph. For many large systems of interest, classical iterative
solvers (such as conjugate gradient and lanczos) converge with
prohibitively slow rates, due to ill-conditioned matrices and small
spectral gap ratios. Scalable preconditioners for this class of problem
are of considerable research interest. We consider constructing algebraic
multigrid (AMG) preconditioners, which make use of coarse approximation
of eigenspaces to accelerate convergence. We discuss several interesting
properties the eigenvectors of scale-free graphs have that are not
present in mesh-like graphs. For example, these graphs have a large
number of locally supported eigenvectors (LSEVs), eigenvectors that are
each nonzero only in a small, connected region. We discuss how such
properties can be used to coarsen graphs with potentially high spectral
accuracy and ultimately improve the scalability of preconditioned
iterative methods.


\end{document}
